\documentclass{situationist-preprint}

\title{Stein's Argument Reconstructed}
\msid{PHE-12-b}
\msfamily{PHE-12}
\msdiscipline{Phenomenology}
\mstype{Companion manuscript}
\mskeywords{Edith Stein; empathy; Einf\"uhlung; non-primordial givenness; Lipps; other minds}

\begin{document}
\maketitle

\begin{abstract}
Edith Stein's dissertation of 1917 contains an argument that empathy is a sui generis act in which another's experience is given directly but non-primordially. The argument is usually reported as a description. We reconstruct it as an argument by elimination over the forms of re-presentation, supplemented by a criticism of the rival theories of Lipps and of inference by analogy. The reconstruction reaches the first conjunct of Theorem~3.3 of PHE-12-a by a route that uses Stein's comparison of empathy with memory, expectation and fantasy, and it shows why Stein held that errors of empathy are corrected by further empathy, which supports the second conjunct.
\end{abstract}

\section{The comparison with re-presentation}

Stein begins from the acts in which something is given without being originally present: memory, expectation and fantasy. In each, the act is primordial, since I live it now, while its content is not: the remembered joy is given as past, the expected one as future, the imagined one as as-if \citep[ch.~2]{Stein1917}. Empathy resembles these acts in that its content is not originally present to me, and differs from all of them. The reconstruction runs as follows.

\begin{assumption}\label{a:forms}
If an experience $e$ is given to me non-primordially, then $e$ is given either as my past experience (memory), as my future experience (expectation), as an experience no one lives (fantasy), or as an experience lived now by another.
\end{assumption}

\begin{assumption}\label{a:present}
In empathy, the other's grief is given as being lived now; it is neither past nor future, and it is not given as merely imagined, since it can surprise and correct me.
\end{assumption}

\begin{assumption}\label{a:notmine}
In empathy, the grief is not given as mine: I need not be grieving, and the grief is located where the other's lived body is.
\end{assumption}

\begin{proposition}\label{p:stein}
In empathy, an experience is given non-primordially as lived now by another.
\end{proposition}

\begin{proof}
The grief is not lived by me now (Assumption~\ref{a:notmine}), so it is given non-primordially. By Assumption~\ref{a:present} it is not given as past, future or merely imagined; by Assumption~\ref{a:notmine} it is not given as my own past or future. Assumption~\ref{a:forms} leaves only the last alternative.
\end{proof}

Proposition~\ref{p:stein} is the non-primordiality half of Stein's thesis, and it agrees with Lemma~3.2 of PHE-12-a. It is reached without the argument used there, which proceeds from the impossibility of another's state being given as mine; here non-primordiality follows from a classification of the acts of re-presentation and the location of empathy among them.

\section{The rival theories}

Directness is established negatively. Lipps held that empathy is an inner imitation in which I feel the other's state, so that, in the full case, I am one with the other \citep{Lipps1903}. Stein replies that the result of such a fusion would be my experience, given primordially, and so not the other's; Lipps's account describes contagion, not empathy \citep[ch.~2]{Stein1917}; \citet{Zahavi2014} presses the same distinction against current simulation theories. The inference from analogy yields a belief that the other is grieving, reached from premises about my own correlations, and not a givenness of the grief; Stein, like \citet{Scheler1923}, observes that the experience of the other's grief does not have the structure of a conclusion. The remaining candidate, association, would at most reproduce my own past states. None of the rivals presents another's present experience as another's. If Proposition~\ref{p:stein} is a correct description, empathy is therefore not reducible to any of them, and its object is given without an intervening inference. This is the directness of Lemma~3.1 of PHE-12-a, reached without Principle~P1 (expressive unity), on which that lemma relies.

Stein distinguishes three levels in the act: the emergence of the other's experience before me, its fulfilling explication, in which I follow it as if from within, and its comprehensive objectification, in which it is again before me, now explicated \citep[ch.~2]{Stein1917}. The structure explains how an act can be direct without being primordial: the middle level is the closest I come to the other's experience, and even there I do not live it as mine.

\begin{objection}
An argument by elimination is only as good as its list of alternatives. Assumption~\ref{a:forms} may be incomplete.
\end{objection}

\begin{reply}
The list is generated by two variables, whose experience the content is given as and when it is given as lived, and the alternatives exhaust their combinations, with fantasy covering the case in which the content is given as lived by no one. A further form would require a content given as lived by no one now and yet not imagined, which describes nothing.
\end{reply}

\section{Knowledge and correction}

Stein holds that empathy can deceive, and that its deceptions are discovered and corrected by further empathy, together with what I know of the other from other sources, not by an independent faculty \citep[ch.~3]{Stein1917}. This is the structure Principle~P2 of PHE-12-a requires: a source of knowledge whose reliability in normal conditions can be checked within its own exercise, as perception is checked by further perception. Theorem~3.3 of PHE-12-a, on which non-primordial givenness suffices for knowledge, is thus Stein's own position, and the reconstruction shows that it does not depend on assimilating empathy to perception.

\begin{thebibliography}{9}\small
\bibitem[Lipps(1903)]{Lipps1903} Lipps, T. (1903). \emph{\"Asthetik: Psychologie des Sch\"onen und der Kunst. Erster Teil}. Leopold Voss.
\bibitem[Scheler(1923)]{Scheler1923} Scheler, M. (1923). \emph{Wesen und Formen der Sympathie}. Friedrich Cohen.
\bibitem[Stein(1917)]{Stein1917} Stein, E. (1917). \emph{Zum Problem der Einf\"uhlung}. Buchdruckerei des Waisenhauses. Trans. W. Stein, \emph{On the Problem of Empathy}, ICS Publications, 1989.
\bibitem[Zahavi(2014)]{Zahavi2014} Zahavi, D. (2014). \emph{Self and Other: Exploring Subjectivity, Empathy, and Shame}. Oxford University Press.
\end{thebibliography}

\end{document}
